\documentclass[11pt,a4paper]{article}
\usepackage[T1]{fontenc}
\usepackage{lmodern}
\usepackage[margin=2.5cm]{geometry}
\usepackage{setspace}
\usepackage{lineno}
\usepackage{amsmath}
\usepackage{graphicx}
\usepackage{booktabs}
\usepackage[tableposition=top]{caption}
\usepackage{authblk}
\usepackage[numbers,sort&compress,super]{natbib}
\usepackage[hidelinks]{hyperref}

% Most medical journals ask for double spacing and line numbers when submitting.
\doublespacing
\linenumbers

\title{Title of the Study: Accuracy of an Index Test for a Target Condition}
\author[1]{First Author}
\author[1,2]{Second Author}
\author[2]{Third Author}
\affil[1]{Department of Medicine, University Hospital, City, Country}
\affil[2]{School of Public Health, University, City, Country}
\date{}

\begin{document}

\maketitle

\noindent\textbf{Corresponding author:} First Author, Department of Medicine,
University Hospital, City, Country; first.author@example.org

\begin{abstract}
\noindent\textbf{Objective:} State which test is studied for which condition.\\
\textbf{Methods:} Describe the participants, the index test and the reference standard.\\
\textbf{Results:} Give sensitivity and specificity with 95\% confidence intervals.\\
\textbf{Conclusions:} Say what the accuracy means for practice.
\end{abstract}

\noindent\textbf{Keywords:} diagnostic accuracy; sensitivity and specificity; ROC curve

\section{Introduction}
Describe the clinical role of the test. The study is reported following STARD.\cite{bossuyt2015stard}

\section{Methods}
\subsection{Participants}
Give the eligibility criteria and how participants were identified.

\subsection{Test methods}
Describe the index test and the reference standard, and who read them.

\subsection{Analysis}
Sensitivity and specificity are
\begin{equation}
  \mathrm{Se} = \frac{\mathrm{TP}}{\mathrm{TP} + \mathrm{FN}}, \qquad
  \mathrm{Sp} = \frac{\mathrm{TN}}{\mathrm{TN} + \mathrm{FP}},
\end{equation}
and the area under the ROC curve summarises discrimination.\cite{hanley1982}

\section{Results}
Table~\ref{tab:accuracy} cross-tabulates the index test with the reference standard.

\begin{table}[htbp]
  \centering
  \caption{Index test results against the reference standard.}
  \label{tab:accuracy}
  \begin{tabular}{lrr}
    \toprule
     & Condition present & Condition absent \\
    \midrule
    Index test positive & 90 & 15 \\
    Index test negative & 10 & 185 \\
    \bottomrule
  \end{tabular}
\end{table}

\section{Discussion}
Discuss the clinical applicability and the limitations.

\section*{Declarations}
\noindent\textbf{Ethics approval:} Name the committee and the approval number.\\
\textbf{Funding:} Name the funders and their role.\\
\textbf{Competing interests:} Declare any, or state that there are none.\\
\textbf{Data availability:} Say where the data can be found.\\
\textbf{Author contributions:} Say who did what.

\bibliographystyle{vancouver}
\bibliography{references}

\end{document}
